Deux petites sphères métalliques identiques portent la même charge électrique. Placées à \qty{\v{r}}{cm} l'une de l'autre, elles se repoussent avec une force de \qty{\v{F}}{N}. La constante de Coulomb est $k = \qty{8.99e9}{N.m^2/C^2}$ et la charge élémentaire vaut \qty{1.602e-19}{C}.

\begin{enumerate}[label=\alph*)]
    \item Déterminez la charge portée par chaque sphère.
    \item Si les sphères sont chargées négativement, combien d'électrons en surplus chacune porte-t-elle ?
    \item Que devient la force si l'on triple la distance entre les sphères ?
\end{enumerate}
